Both $I,H$ lie on the same side of $BC$, so cyclicity gives $\angle BIC=\angle BHC$. Thus $90^\circ+\dfrac{A}{2}=180^\circ-A$, or $A=60^\circ$. Euler's line gives $OH=3OG=240$. The circle through $B,C,H$ is the reflection of the circumcircle across $BC$, since the reflection of the orthocentre across a side lies on the circumcircle (its angle at the reflected point is supplementary to $A$). Therefore its radius is $R=312$, and its centre $J$ is the reflection of $O$. As the distance from $O$ to $BC$ is $R\cos60^\circ=156$, $OJ=312=R$, so $O$ lies on $\omega$. The diameter theorem makes $\angle HOH^{\prime}=90^\circ$, giving $OH^{\prime}=\sqrt{624^2-240^2}=576$. Use coordinates $O=(0,0)$, $J=(0,-R)$ and $BC:y=-\dfrac{R}{2}$. If $H^{\prime}=(u,v)$, its circle equation gives $576^2=u^2+v^2=-2Rv$. Since $P$ lies on ray $OH^{\prime}$ at height $-\dfrac{R}{2}$, similarity along that ray gives $OP=576\cdot\dfrac{\dfrac{R}{2}}{\dfrac{576^2}{2R}}=\dfrac{312^2}{576}=\boxed{169}$.
\begin{asy}
import graph;import geometry;usepackage("amsmath");size(7cm);
real R=312,ay=2856/13;pair O=(0,0),J=(0,-312),A=(sqrt(R*R-ay*ay),ay),B=(-156sqrt(3),-156),C=(156sqrt(3),-156),H=A+(0,-312),Hp=-A+(0,-312),P=Hp*(-156/Hp.y),G=H/3;real a=abs(B-C),b=abs(A-C),c=abs(A-B);pair I=(a*A+b*B+c*C)/(a+b+c);draw(A--B--C--cycle);draw(circle(J,R),gray);draw(H--Hp^^O--Hp,dashed);dot(O);dot(I);dot(H);dot(G);label("$A$",A,(0,1));label("$B$",B,(-1,0));label("$C$",C,(1,0));label("$O$",O,(-1,1));label("$I$",I,(-1,-1));label("$H$",H,(1,0));label("$G$",G,(1,1));label("$H^{\prime}$",Hp,(0,-1));label("$P$",P,(-1,0));
\end{asy}

